\phantomsection
\addcontentsline{toc}{chapter}{\bf ABSTRACT}

\begin{center}
    \large{\bf ABSTRACT}
\end{center}

Longitudinal comparison of medical images is a cornerstone of clinical practice, essential for tracking disease progression and evaluating treatment response. To automate this process, the field of Difference Visual Question Answering (DiffVQA) has emerged, yet current state-of-the-art models suffer from critical limitations. They are computationally expensive, processing thousands of dense visual features; they lack explicit mechanisms to reason about the direction of clinical change; and they are vulnerable to hallucinating findings not present in the visual evidence. To address these challenges, this thesis introduces DRIFT-QA, a novel, lightweight architecture designed for a new paradigm of efficiency, explicit reasoning, and robustness. Our model features three key innovations: a Directional Residual Stack (DRS) to explicitly model whether a finding has worsened or improved; a Question-Guided Difference Tokenizer (QDT) that uses cross-attention to distill the visual difference into a sparse set of just a few relevant tokens; and a Counterfactual Regularizer (CFA-Reg) that trains the model to be robust against hallucinations. The author's experiments show that DRIFT-QA achieves state-of-the-art performance on the MIMIC-Diff-VQA benchmark, outperforming larger, brute-force models while using a fraction of the computational resources. The significance of this work is threefold: it introduces a new, efficient paradigm for DiffVQA, contributes several novel and independently valuable architectural components, and establishes a new benchmark for creating models that are not only accurate but also practical and trustworthy for clinical support.
